\chapter{Le socle mécaniste : fréquence et sévérité}
\label{chap:socle}

\noindent Ces deux briques (loi de sévérité GPD par pilier, fréquence surdispersée pilotée par
un facteur commun) sont reprises telles quelles : elles alimentent l'agrégation, qui est le
seul élément que la cascade remplace.

% MIGRÉ depuis memoire/main.tex (plages 787-932, 1492-1667).
% Hiérarchie descendue d'un cran, labels préfixés "soc:".
% RELECTURE FAITE (août 2026). Vérifié : aucune phrase de ce chapitre n'annonce la
% copule de Gumbel ni la variable latente comme LE modèle du mémoire ; la copule
% n'apparaît qu'au chapitre 13 comme axe de comparaison de la bande de modèle.
% Deux ambiguïtés de l'ancien modèle ont été levées ici :
%   (a) l'encadré "deux grilles coexistent" articule vecteur d'attaque et pilier ;
%   (b) l'encadré sur lambda_ref dit explicitement que les 341 de la PRC sont une clé
%       de répartition et n'alimentent PAS le moteur (l'entrée est lambda = 21,6).
% La charge systémique empirique a été corrigée sur la sortie réelle du script 08b
% (a = 0,122 [0,054 ; 0,223] ; phi = 2,24), l'ancien "a ~ 0,30 / phi ~ 3" était faux.
%--------------------------------------------------------------

Là où le Théorème Central Limite s'intéresse au comportement asymptotique des sommes de variables aléatoires, la Théorie des Valeurs Extrêmes (EVT) étudie le comportement asymptotique de leurs maxima.

Soit une séquence $(X_1, X_2, \dots, X_n)$ de variables aléatoires indépendantes et identiquement distribuées (i.i.d.) de fonction de répartition $F$, représentant la sévérité des cyberattaques. L'approche initiale de l'EVT, dite des \textit{Block Maxima}, s'intéresse au maximum sur des blocs de taille $n$ (par exemple, le sinistre cyber maximal annuel) :
\begin{equation}
M_n = \max(X_1, X_2, \dots, X_n)
\end{equation}

\section{Les deux théorèmes qui autorisent l'extrapolation de queue}
% HARNAIS-HORS-SECTION: enonces de theoremes. Les nombres y sont des constantes d'enonce
%   (indices de domaine d'attraction, bornes, dates de publication), aucun n'est calcule
%   par le pipeline ; les valeurs empiriques correspondantes sont verifiees dans les
%   sections de calibration et de validation du meme chapitre.

\paragraph{Théorème de Fisher-Tippett-Gnedenko (1928, 1943).} S'il existe des suites de constantes de normalisation $a_n > 0$ et $b_n \in \mathbb{R}$, et une distribution non dégénérée $H$ telles que :
\begin{equation}
\lim_{n \to \infty} \mathbb{P}\left( \frac{M_n - b_n}{a_n} \le x \right) = \lim_{n \to \infty} F^n(a_n x + b_n) = H(x)
\end{equation}
Alors, $H$ appartient nécessairement à la famille des distributions des extrêmes généralisées (GEV - \textit{Generalized Extreme Value}), dont la fonction de répartition s'écrit :
\begin{equation}
H_{\xi}(x) =
\begin{cases}
\exp\left(-\left(1 + \xi x \right)^{-1/\xi}\right) & \text{si } \xi \neq 0 \\
\exp\left(-\exp(-x)\right) & \text{si } \xi = 0
\end{cases}
\end{equation}
définie sur l'ensemble $\{x : 1 + \xi x > 0\}$. Le paramètre $\xi$ (indice de queue) définit trois grands domaines d'attraction :
\begin{itemize}
    \item \textbf{$\xi > 0$ (Loi de Fréchet) :} La distribution présente une queue lourde à décroissance polynomiale. C'est le régime qui décrit les sinistres cyber les plus lourds.
    \item \textbf{$\xi = 0$ (Loi de Gumbel) :} La distribution présente une queue légère à décroissance exponentielle (domaine d'attraction des lois Normale, Lognormale, Gamma).
    \item \textbf{$\xi < 0$ (Loi de Weibull) :} La distribution a un support fini supérieur (queue courte).
\end{itemize}

\noindent On peut énoncer formellement ce résultat fondateur, qui joue pour les maxima le rôle du théorème central limite pour les sommes.

\begin{theorem}[Fisher--Tippett--Gnedenko]\label{soc:th:ftg}
Soit $(X_i)_{i\ge1}$ i.i.d.\ de fonction de répartition $F$ et $M_n=\max(X_1,\dots,X_n)$. S'il existe des suites $a_n>0$, $b_n\in\R$ et une loi non dégénérée $H$ telles que $(M_n-b_n)/a_n \xrightarrow{d} H$, alors $H$ est, à un changement d'échelle et de position près, une loi GEV $H_\xi$. Le signe de $\xi$ détermine le domaine d'attraction : Fréchet ($\xi>0$, queue lourde), Gumbel ($\xi=0$) ou Weibull ($\xi<0$, support borné).
\end{theorem}

\begin{proof}[Éléments]
La loi limite $H$ est nécessairement \emph{max-stable} : pour tout $n$, $H^n(a_nx+b_n)=H(x)$, car $M_{nk}$ est le maximum de $k$ blocs de taille $n$. L'équation fonctionnelle de max-stabilité, résolue via la théorie des fonctions à variation régulière, n'admet (à normalisation près) que la famille à un paramètre $H_\xi$. La correspondance queue de $F$ $\leftrightarrow$ signe de $\xi$ s'obtient en caractérisant les $F$ dont la fonction de survie est à variation régulière d'indice $-1/\xi$ (cas Fréchet). Le détail figure en Annexe~\ref{soc:ann:ftg}.
\end{proof}

%--------------------------------------------------------------

%--------------------------------------------------------------

L'approche \textit{Block Maxima} (GEV) est rigoureuse mais coûteuse en information. En ne
retenant que le maximum de chaque bloc, elle écarte le deuxième plus grand sinistre d'une année
extrême, lequel peut dépasser le maximum d'une année calme. Sur une donnée cyber aussi peu
fournie, ce rejet d'observations fragiliserait l'ajustement.

Pour contourner ce problème, ce mémoire privilégie l'approche \textit{Peaks-Over-Threshold} (POT). Au lieu de ne garder que les maxima, cette méthode exploite toutes les observations excédant un seuil élevé $u$. On s'intéresse alors à la distribution des excès $Y = X - u$.

La fonction de répartition conditionnelle des excès est donnée par :
\begin{equation}
F_u(y) = \mathbb{P}(X - u \le y \mid X > u) = \frac{F(u+y) - F(u)}{1 - F(u)}
\end{equation}

\paragraph{Théorème de Balkema-de Haan-Pickands (1974).} Ce théorème établit le pont mathématique entre la GEV et l'approche POT. Pour une large classe de distributions sous-jacentes $F$ (celles appartenant au domaine d'attraction d'une loi GEV $H_\xi$), il existe une fonction mesurable positive $\sigma(u)$ telle que, lorsque le seuil $u$ tend vers la limite supérieure du support de $F$ ($x_F$), la distribution des excès converge asymptotiquement vers une distribution de Pareto Généralisée (GPD) :
\begin{equation}
\lim_{u \to x_F} \sup_{0 \le y < x_F - u} \left| F_u(y) - G_{\xi, \sigma(u)}(y) \right| = 0
\end{equation}

Où la fonction de répartition de la GPD s'écrit :
\begin{equation}
G_{\xi, \sigma}(y) =
\begin{cases}
1 - \left(1 + \frac{\xi y}{\sigma}\right)^{-1/\xi} & \text{si } \xi \neq 0 \\
1 - \exp\left(-\frac{y}{\sigma}\right) & \text{si } \xi = 0
\end{cases}
\end{equation}
avec $y \ge 0$ si $\xi \ge 0$, et $0 \le y \le -\sigma/\xi$ si $\xi < 0$. Le paramètre de forme $\xi$ de la GPD est strictement le même que celui de la loi GEV correspondante. Dans le contexte du risque cyber, un indice de queue $\xi > 0$ traduit une distribution à queue lourde de type Fréchet, impliquant des moments potentiellement infinis (l'espérance est infinie si $\xi \ge 1$, la variance si $\xi \ge 0.5$). Le résultat s'énonce formellement ainsi.

\begin{theorem}[Pickands--Balkema--de Haan]\label{soc:th:balkema}
Soit $F$ dans le domaine d'attraction d'une GEV $H_\xi$ (Théorème~\ref{soc:th:ftg}). Alors il existe une fonction $\sigma(u)>0$ telle que la loi des excès converge uniformément vers une GPD lorsque le seuil approche la borne supérieure $x_F$ du support :
\begin{equation}
\lim_{u\to x_F}\ \sup_{0\le y< x_F-u}\bigl|F_u(y)-G_{\xi,\sigma(u)}(y)\bigr|=0,
\end{equation}
où $F_u(y)=\Prob(X-u\le y\mid X>u)$ et $G_{\xi,\sigma}$ est la fonction de répartition de la GPD. L'indice de queue $\xi$ est identique à celui de la GEV limite.
\end{theorem}

\begin{proof}[Esquisse]
L'appartenance de $F$ au domaine d'attraction de $H_\xi$ équivaut à une condition de variation régulière sur la fonction de survie $\bar F$. On montre que $\bar F(u+y\,a(u))/\bar F(u)\to(1+\xi y)^{-1/\xi}$ quand $u\to x_F$, pour une fonction d'échelle $a(u)$ adéquate ; ce membre de droite est exactement la fonction de survie GPD. La convergence uniforme découle de la monotonie des fonctions de répartition (théorème de Pólya). Détail et lien avec le Théorème~\ref{soc:th:ftg} en Annexe~\ref{soc:ann:balkema}.
\end{proof}

%--------------------------------------------------------------

%--------------------------------------------------------------

L'estimation de l'indice de queue $\xi$ est un enjeu critique, car une infime variation de ce paramètre modifie massivement le Capital de Solvabilité Requis (SCR). Deux approches principales s'affrontent dans la littérature.

\section{Estimer l'indice de queue, et savoir ce que vaut l'estimation}

% SOURCES-SCRIPTS: 47 63 67

\paragraph{L'estimateur de Hill.}
Pour $\xi > 0$, l'estimateur non paramétrique de Hill (1975) estime l'indice de queue à partir des $k$ plus grandes statistiques d'ordre, sans passer par l'ajustement d'une loi paramétrique. Asymptotiquement convergent et sans biais lorsque la loi appartient au domaine d'attraction de Fréchet, il présente en échantillon fini un biais majeur : si la composante à variation lente de la queue ne converge pas assez vite, $\hat{\xi}_{\text{Hill}}$ devient instable.

\paragraph{Le Maximum de Vraisemblance (MLE).}
Pour pallier le biais de l'estimateur de Hill hors du régime asymptotique strict, ce mémoire privilégie l'estimation par Maximum de Vraisemblance (MLE) sur les excès de la GPD. La log-vraisemblance pour un échantillon de $N_u$ excès $y_i$ au-delà du seuil $u$ s'écrit (pour $\xi \neq 0$) :
\begin{equation}
\ell(\xi, \sigma) = -N_u \ln(\sigma) - \left(1 + \frac{1}{\xi}\right) \sum_{i=1}^{N_u} \ln\left(1 + \frac{\xi y_i}{\sigma}\right)
\label{soc:eq:loglik-gpd}
\end{equation}
La maximisation numérique de cette fonction fournit des estimateurs $(\hat{\xi}_{MLE}, \hat{\sigma}_{MLE})$ qui, bien que nécessitant des méthodes d'optimisation (algorithme de Nelder-Mead), offrent une robustesse supérieure face à la variance d'échantillonnage des bases de sinistres cyber comme OpRisk ou PRC.

\paragraph{Loi asymptotique et information de Fisher.} La qualité d'un intervalle de confiance sur $\xi$ conditionne toute la lecture du capital. Le fondement théorique en est donné ici ; il sert de contrôle indépendant du bootstrap employé au chapitre~\ref{chap:resultats}.

\begin{proposition}[Normalité asymptotique de l'EMV de la GPD]\label{soc:prop:fisher}
Pour $\xi > -\tfrac{1}{2}$, l'estimateur du maximum de vraisemblance $(\hat\xi,\hat\sigma)$ de la GPD est consistant et asymptotiquement normal :
\begin{equation}
\sqrt{N_u}\,\begin{pmatrix}\hat\xi-\xi\\[2pt]\hat\sigma-\sigma\end{pmatrix}
\;\xrightarrow{\;d\;}\;
\mathcal{N}\!\left(\mathbf{0},\; (1+\xi)\begin{pmatrix}1+\xi & -\sigma\\[2pt] -\sigma & 2\sigma^2\end{pmatrix}\right).
\end{equation}
En particulier, la variance asymptotique de l'indice de queue est
$\Var(\hat\xi)\simeq (1+\xi)^2/N_u$.
\end{proposition}

\begin{proof}[Esquisse]
La matrice de covariance limite est l'inverse de l'information de Fisher $\mathcal{I}(\xi,\sigma)$, obtenue en prenant l'espérance de l'opposé de la hessienne de $\ell$. Le détail du calcul de $\mathcal{I}$ (dérivées secondes de \eqref{soc:eq:loglik-gpd} et intégration contre la densité GPD, valide pour $\xi>-\tfrac12$ où l'information est finie) est reporté en Annexe~\ref{soc:ann:fisher}. Le résultat classique remonte à \citet{smith1987}.
\end{proof}

\noindent Appliquée à la calibration OpRisk ($\hat\xi=0{,}595$, $N_u=91$ excès), la Proposition~\ref{soc:prop:fisher} donne un écart-type asymptotique $\widehat{\mathrm{sd}}(\hat\xi)=(1+\hat\xi)/\sqrt{N_u}=0{,}167$, soit un intervalle de confiance à 90\,\% par la delta-méthode de $[0{,}320\,;\,0{,}870]$. Cet intervalle recoupe presque exactement l'intervalle bootstrap $[0{,}304\,;\,0{,}831]$ obtenu indépendamment au \S\ref{soc:sec:severite} : la concordance des deux approches, l'une paramétrique et asymptotique, l'autre par rééchantillonnage, renforce la fiabilité de l'incertitude reportée sur~$\xi$.

\begin{proposition}[Loi asymptotique de l'estimateur de Hill]\label{soc:prop:hill}
Si $F$ est à variation régulière d'indice $-1/\xi$ (domaine de Fréchet), alors pour une suite $k=k(n)\to\infty$ avec $k/n\to0$,
\begin{equation}
\sqrt{k}\,\bigl(\hat\xi_{Hill}(k)-\xi\bigr)\;\xrightarrow{\;d\;}\;\mathcal{N}(0,\;\xi^2).
\end{equation}
\end{proposition}

\noindent L'écart-type $\xi/\sqrt{k}$ décroît en $\sqrt{k}$ : d'où l'arbitrage biais--variance qui gouverne le choix de $k$ dans le \emph{Hill plot}, et qui explique l'instabilité signalée plus haut lorsque $k$ est pris trop grand (contamination par le corps de la loi) ou trop petit (variance explosive). La preuve, fondée sur la représentation de Rényi des espacements exponentiels, est esquissée en Annexe~\ref{soc:ann:hill}.

\paragraph{Quantile de sévérité : formules fermées, et ce qu'elles ne donnent pas.} Le passage des paramètres $(\xi,\sigma)$ au quantile d'\emph{un} sinistre est explicite. Il faut dire tout de suite ce qu'il n'est pas : ce quantile n'est pas un besoin de capital, et l'on n'y accède au capital qu'après agrégation annuelle de la fréquence et de la sévérité, puis propagation dans la cascade (chapitre~\ref{chap:cascade}). Notons $\zeta_u=\Prob(X>u)$ la probabilité de dépassement du seuil. Les formes qui suivent sont
classiques \citep{mcneil2015}.

\begin{proposition}[Quantile et moyenne de queue d'une sévérité GPD]\label{soc:prop:varestvar}
Sous le modèle POT, pour $\alpha$ tel que $\qsev_\alpha(X)>u$ et $\xi\in(0,1)$ :
\begin{align}
\qsev_\alpha(X) &= \boxed{\,u + \frac{\sigma}{\xi}\left[\left(\frac{1-\alpha}{\zeta_u}\right)^{-\xi}-1\right]\,},\label{soc:eq:var-gpd}\\[4pt]
\qbarsev_\alpha(X) = \E[X\mid X>\qsev_\alpha(X)] &= \boxed{\,\frac{\qsev_\alpha(X)}{1-\xi}+\frac{\sigma-\xi u}{1-\xi}\,}.\label{soc:eq:tvar-gpd}
\end{align}
Lorsque $\xi\ge 1$, l'espérance $\E[X]$ est infinie et $\qbarsev_\alpha$ n'est pas définie.
\end{proposition}

\begin{proof}
Pour $x>u$, la propriété POT donne $\Prob(X>x)=\zeta_u\bigl(1+\xi(x-u)/\sigma\bigr)^{-1/\xi}$. En posant ce membre égal à $1-\alpha$ et en inversant, on obtient \eqref{soc:eq:var-gpd}. Pour \eqref{soc:eq:tvar-gpd}, la stabilité de la GPD par franchissement de seuil implique que l'excès au-delà de $\qsev_\alpha(X)$ suit encore une GPD de paramètre $\xi$ et d'échelle $\sigma+\xi(\qsev_\alpha(X)-u)$, d'espérance $\bigl(\sigma+\xi(\qsev_\alpha(X)-u)\bigr)/(1-\xi)$ pour $\xi<1$ ; en ajoutant $\qsev_\alpha(X)$ et en réarrangeant, on retrouve \eqref{soc:eq:tvar-gpd}. La divergence pour $\xi\ge1$ découle de $\E[X]=\int_u^\infty \Prob(X>x)\,\dd x=+\infty$ dès que $1/\xi\le1$. Détail en Annexe~\ref{soc:ann:varestvar}.
\end{proof}

\noindent Numériquement, la calibration OpRisk ($u=20{,}03$, $\sigma=57{,}97$, $\xi=0{,}595$, $\zeta_u=0{,}151$) donne par \eqref{soc:eq:var-gpd} un quantile de sévérité $\qsev_{99{,}5\%}(X)=663{,}0$~M\euro\ et par \eqref{soc:eq:tvar-gpd} une moyenne de queue $\qbarsev_{99{,}5\%}(X)=1\,752{,}4$~M\euro.

\begin{attention}
\textbf{Ces deux nombres ne sont pas des capitaux, et le mémoire ne les compare jamais à un SCR.}
Ils décrivent \emph{un} sinistre. La mesure de capital est $\VaR_{99{,}5\,\%}(L)$, quantile de la
charge \emph{annuelle agrégée}, et elle vaut plusieurs milliers de millions d'euros à l'échelle
sectorielle : l'écart avec les $663$~M\euro\ ci-dessus ne traduit aucune incohérence, il traduit
l'agrégation. À $\lambda\simeq21{,}6$ sinistres matériels par an, atteindre le quantile annuel à
$99{,}5\,\%$ demande un quantile \emph{par sinistre} de l'ordre de $99{,}98\,\%$, bien plus haut
dans la queue. Le pont entre les deux échelles est chiffré au chapitre~\ref{chap:resultats}.
% SOURCES-SCRIPTS: 67
Réserver $\VaR$ et $\TVaR$ à $L$, et $\qsev$ et $\qbarsev$ à
$X$, n'est donc pas une coquetterie de notation : c'est ce qui empêche de lire une différence
d'échelle comme une contradiction.
\end{attention}

\noindent La Proposition~\ref{soc:prop:varestvar} éclaire par ailleurs le contraste entre sources : la calibration PRC affiche $\hat\xi=1{,}033>1$, ce qui place la sévérité en régime d'espérance infinie : la moyenne de queue, et donc tout Expected Shortfall bâti sur elle, y est mathématiquement indéfinie. C'est ce qui justifie \emph{a posteriori} le recours à un plafond de sévérité (capacité de réassurance) pour rendre le capital PRC calculable, et c'est aussi le premier argument contre l'usage d'une mesure de queue comme \emph{niveau de couverture} : sur l'une des deux sources, cette mesure n'existe pas.

%--------------------------------------------------------------


%--------------------------------------------------------------

\begin{attention}
\textbf{Deux grilles coexistent dans ce chapitre, et elles ne sont pas concurrentes.} La
première décrit \emph{comment} un incident entre : par vecteur d'attaque (hameçonnage,
exploitation de vulnérabilité, chaîne d'approvisionnement). La seconde décrit \emph{quel
contrôle} a cédé : par pilier DORA. Un vecteur n'est pas un pilier, et il serait faux
de les apparier terme à terme. Leur articulation est de type \emph{entrée puis propagation} :
le vecteur gouverne l'amorce, c'est-à-dire la fréquence d'arrivée d'un sinistre et le
pilier sur lequel il tombe ; la cascade du chapitre~\ref{chap:cascade} gouverne ensuite la
\emph{propagation} entre piliers. Le présent chapitre calibre donc l'amorce et rien d'autre.
La grille par vecteur ne réapparaît plus après lui, et aucun résultat de capital ne dépend de
son appariement à un pilier.
\end{attention}

\section{La brique de fréquence}

\paragraph{Décomposition par vecteur d'attaque.} Une fréquence de référence
$\lambda_{ref}$, lue sur le périmètre financier de la base PRC
($\lambda_{ref} = 341$ notifications/an), est décomposée en
% SOURCES-SCRIPTS: 62
sous-fréquences par vecteur d'attaque,
pondérées par la structure observée dans la base Hackmageddon :
\begin{equation}
\lambda_{ref} = \sum_{v} \lambda_v,
\qquad
\lambda_v = \lambda_{ref} \times \pi_v,
\end{equation}
où $\pi_v$ est la part du vecteur $v$ dans la structure Hackmageddon
(38{,}8~\% phishing, 33{,}8~\% exploitation de vulnérabilités, 15{,}8~\%
supply chain, 6{,}3~\% identifiants compromis, 5{,}3~\% autres). Cette
décomposition transforme un constat structurel en levier de modélisation :
un défaut de conformité sur une exigence donnée n'agit que sur la
sous-fréquence du vecteur qu'elle est censée mitiger, pas sur l'ensemble de
$\lambda_{ref}$.

\begin{attention}
\textbf{$\lambda_{ref}$ est une clé de répartition, pas le niveau du modèle.} Les $341$ de la
PRC comptent des \emph{notifications de violation de données} déclarées par les organisations
financières de la base, sur un périmètre américain et sans seuil de matérialité. Le périmètre
importe : la base entière en compte $2\,150$ par an, six fois plus, et confondre les deux
rendrait le nombre incompréhensible à qui le recalcule. Ce nombre sert exclusivement à
répartir la sinistralité entre
vecteurs, c'est-à-dire à fixer les $\pi_v$ et à faire jouer les multiplicateurs de conformité.
\textbf{Il n'alimente pas le moteur de cascade.} L'entrée du moteur est la fréquence
d'incidents \textsc{tic} \emph{matériels} lue sur OpRisk, $\lambda \simeq 21{,}6$ par an au
niveau du secteur financier mondial, dont la table du chapitre~\ref{chap:donnees} donne les
trois échelles de lecture. Tout niveau de capital du chapitre~\ref{chap:resultats} repose sur
$\lambda$, jamais sur $\lambda_{ref}$.
\end{attention}

\paragraph{Choix de la loi binomiale négative.} La fréquence agrégée
$\lambda_{S_x} = \sum_v \lambda_{ref}\,\pi_v\, m_{v,S_x}$ obtenue sous un
scénario $S_x$ n'est pas simulée par une loi de Poisson, mais par une loi
binomiale négative, pour tenir compte de la sur-dispersion observée dans les
données ($\sigma^2/\mu > 1$). Le facteur de sur-dispersion
$\phi = \sigma^2/\mu = 9{,}2$ est traité comme une propriété structurelle de
la sinistralité, supposée indépendante du niveau de fréquence : c'est donc
$\phi$ qui est maintenu constant d'un scénario à l'autre, et non le
paramètre $r$ de la binomiale négative, qui est au contraire recalculé sous
chaque scénario :
\begin{equation}
r_{S_x} = \frac{\lambda_{S_x}}{\phi - 1},
\qquad
p_{S_x} = \frac{r_{S_x}}{r_{S_x} + \lambda_{S_x}}.
\end{equation}
Augmenter $\lambda$ via les multiplicateurs DORA déplace ainsi la fréquence
moyenne tout en préservant la \emph{forme} de sur-dispersion de la
sinistralité, hypothèse de modélisation explicite, dont la limite est
que le comportement de la sur-dispersion en régime de forte dégradation
n'est pas lui-même observé, faute d'historique de crise DORA.

\paragraph{Réconciliation des quatre mesures de surdispersion, et effet sur le capital.}
Avant les trois mesures de forme, un mot sur celle que l'on observe directement, car un lecteur
qui parcourt ce chapitre rencontre \emph{quatre} nombres de surdispersion et pourrait les croire
en concurrence. Ils ne le sont pas, ils vivent à des échelles différentes, et
l'indice de dispersion n'est pas invariant d'échelle : pour un mélange de Poisson il
vaut $1+\lambda/r$, donc il croît avec le niveau de fréquence auquel on le lit. Sur les
$14\,944$ cellules firme-année du panel, où la fréquence moyenne est de l'ordre de quatre
centièmes d'incident par cellule, on mesure $\Var/\E=1{,}19$. Sur la série annuelle sectorielle
détendancée, où $\lambda$ vaut $21{,}6$, on mesure $2{,}24$. Le moteur, lui, retient $9{,}20$,
qui est un choix de prudence posé et non une mesure. Une échelle de lecture, une
agrégation de cellules corrélées, puis une marge prudentielle assumée : les trois écarts
s'expliquent l'un après l'autre, et aucun ne contredit les autres.
Le modèle porte, à trois échelles distinctes, trois mesures de surdispersion qu'il faut
réconcilier pour ne pas les confondre : la forme du mélange Gamma au niveau firme-année
($r=0{,}63$), la charge systémique commune ($a=0{,}60$, mélange lognormal du
chapitre~\ref{chap:multietats}) et le facteur $\varphi=9{,}20$ ci-dessus (mélange Gamma
agrégé, c'est-à-dire la binomiale négative). Un mélange de Poisson a pour indice de dispersion
$\Var/\E$ égal à $1+\lambda/r=\varphi$ (mélange Gamma) ou $1+\lambda(e^{a^2}-1)$ (mélange
lognormal) ; les deux décrivent \emph{la même} surdispersion agrégée dès lors que
$a=\sqrt{\ln\!\bigl(1+(\varphi-1)/\lambda\bigr)}$. À l'échelle du moteur ($\lambda\simeq 21{,}6$
entrées par an), $\varphi=9{,}20$ correspond à $a=0{,}57$, soit à cinq pour cent de la valeur
$0{,}60$ posée par ailleurs : les deux calages agrégés coïncident, le $r=0{,}63$ firme vivant à
une échelle firme-année non comparable directement.

Trois enseignements suivent (figure~\ref{soc:fig:nb}). D'abord, à variance égale la queue
dépend encore de la loi de mélange : le quantile $99{,}5\,\%$ du nombre annuel vaut $74$
(mélange Gamma) contre $82$ (lognormal), soit $+11\,\%$ : appairer la variance n'appaire pas
la queue. Ensuite, le choix de la loi de comptage déplace la VaR mais \emph{non} la moyenne
(l'espérance $\E[N]=\lambda$ est fixe pour les trois lois) : passer du Poisson à la binomiale
négative posée relève le SCR de $+9\,\%$ seulement, la queue restant portée au $99{,}5\,\%$ par
la perte unique dominante ($\xi\simeq 0{,}60$). Enfin, calée sur la série cyber$\times$finance
détrendée, le rapport $\Var/\E$ brut confondant la croissance tendancielle avec le choc
commun, la dispersion de Pearson résiduelle vaut $\varphi=2{,}24$
(IC~$95\,\%$ $[1{,}24\,;5{,}18]$, dix-huit points seulement), soit une charge empirique
$a=0{,}122$ (IC~$95\,\%$ $[0{,}054\,;0{,}223]$). La valeur posée $a=0{,}60$ dépasse donc la
\emph{borne haute} de cet intervalle, et le $\varphi=9{,}20$ retenu au niveau du moteur
dépasse de loin le $2{,}24$ mesuré : l'un comme l'autre sont un choix de prudence sur
la queue de fréquence, non une calibration, à assumer comme tel. La conclusion
\og{}le calage posé est conservateur\fg{} résiste ainsi à l'incertitude d'estimation, ce qui
est la seule chose que dix-huit points autorisent à affirmer.
% SOURCES-SCRIPTS: 47 08b 67 41

\paragraph{La dynamique du facteur systémique.} Le facteur commun $Y_{j,t}$ est posé centré
et réduit. Sa dynamique, écrite au \S\ref{sec:dynamique-Y}, superpose un bruit de fond
brownien et une composante de sauts pour les campagnes de masse. Elle ne déplace aucune
probabilité marginale, donc aucun capital publié, et n'agit que sur la dépendance de queue ; la
forme retenue est la forme à renouvellement, sans mémoire d'un exercice à l'autre.

\begin{figure}[htbp]\centering
\figover{N2_simulation_nb.png}
\caption{La loi de comptage négative binomiale, simulée et réconciliée. (a)~à moyenne égale,
la queue du nombre annuel diffère selon la loi de mélange. (b)~la charge lognormale $a$ et le
facteur $\varphi$ tracent la même surdispersion agrégée ; les deux valeurs posées quasi
coïncident. (c)~le comptage déplace la VaR, pas la moyenne.}
% SOURCES-SCRIPTS: 41
\label{soc:fig:nb}
\end{figure}

\paragraph{De la fréquence par scénario à la fréquence par entité.} Le choix
d'un scénario $S_x \in \{S_0, S_1, S_2\}$ reste, par construction, discret et
appliqué uniformément à l'ensemble du marché. L'état de conformité continu
$C^\star_j$ du chapitre~\ref{chap:multietats} permet de dépasser cette
limite, en substituant à ce choix discret une interpolation continue des multiplicateurs
$m_v$ le long de la latente de conformité,
où $\bar{m}_v^{S_2}$ est le centre de l'intervalle de multiplicateur sourcé
en scénario $S_2$. Une entité au profil de conformité mature conserve ainsi
une fréquence proche de $S_0$ même en cas de choc systémique $Y$
défavorable, tandis qu'une entité en retard de conformité peut, sous un
$\rho_i$ élevé, dépasser le multiplicateur $S_2$ dès l'environnement normal.

%--------------------------------------------------------------

%--------------------------------------------------------------

\section{La brique de sévérité, calibrée sur OpRisk}
\label{soc:sec:severite}

% SOURCES-SCRIPTS: 07 47 63 67

\paragraph{Choix de la loi de Pareto généralisée.} La sévérité des sinistres
cyber présente une signature de queue lourde documentée au chapitre~\ref{chap:donnees} : sur
OpRisk Global, sur la population de la fréquence et de la descente d'échelle, les $10\,\%$ d'incidents les plus
graves concentrent $92{,}3\,\%$
de la perte totale. Cette concentration extrême justifie le recours à la
théorie des valeurs extrêmes, et plus précisément à une modélisation de la
queue de distribution par une loi de Pareto généralisée (GPD) au-delà d'un
seuil $u$, selon l'approche \emph{peaks-over-threshold}.

\paragraph{Calibration OpRisk.} Au seuil du percentile~85
($u = 20{,}0$~M\euro, 91 excès), la calibration donne :
\begin{equation}
\hat{\xi}_{\text{OpRisk}} = 0{,}595 \quad (\text{IC}_{90\%}=[0{,}30\,;\,0{,}83]),
\qquad
\hat{\sigma}_{\text{OpRisk}} = 58{,}0~\text{M\euro},
\qquad
\qsev_{99{,}5\%}(X) = 663~\text{M\euro}.
\end{equation}
Le dernier terme est un quantile de sévérité, c'est-à-dire le montant d'\emph{un}
sinistre, et non un besoin de capital. L'intervalle de confiance bootstrap qui l'entoure,
$[411\,;\,1037]$~M\euro, matérialise à lui seul un facteur~2{,}5 entre les bornes, uniquement
imputable à l'incertitude de paramètre.

\paragraph{Ce que le seuil doit aux diagnostics, et ce qu'il ne leur doit pas.} Le choix de $u$
s'appuie sur deux diagnostics, à condition de ne leur faire dire ni plus ni autre chose que ce
qu'ils disent. Le \emph{diagnostic GPD} (figure~\ref{soc:fig:diagnostic-gpd}) établit un point,
et un seul : le seuil retenu se situe à la frontière où le paramètre de forme quitte la zone
instable ($\xi \geq 1$, où l'espérance n'existe pas) pour un régime de moments finis, tout en
conservant assez d'excès pour estimer ($n = 91 > 30$). C'est un argument de voisinage,
et il ne s'étend pas à toute la plage : en portant le seuil au percentile $90$ puis au percentile
$95$ des données, l'indice descend de $0{,}60$ à $0{,}53$ puis $0{,}38$. Écrire que
$\hat\xi$ est \emph{stable} quand on fait varier le seuil serait donc inexact ; ce qui est
stable, c'est son voisinage, où il ne bouge que de $0{,}5\,\%$. Cette décroissance est mesurée
au \S\ref{soc:sec:validation}, avec ce qu'elle change et ce qu'elle ne change pas. Quant au
\emph{Hill plot} (figure~\ref{soc:fig:hill-ic}), il ne corrobore pas ce niveau par sa valeur,
qui en vaut plus du double : la section suivante établit que cet écart est exactement ce qu'un
$\xi$ de $0{,}60$ produit en échantillon fini, ce qui en fait une corroboration d'un autre
genre, et plus forte que la lecture d'un plateau.

\begin{figure}[htbp]
\centering
\figover{figures/diagnostic_gpd_oprisk.png}
\caption{Diagnostic GPD sur le périmètre cyber$\times$finance d'OpRisk.
\textit{Gauche} : $\hat{\xi}$ en fonction du seuil $u$ ; le paramètre passe sous
$\xi = 1$ (espérance finie) avant le seuil retenu $u = 20$~M\euro. La plage tracée
s'arrête en deçà des seuils les plus élevés du balayage en percentiles, où l'indice
descend plus bas : cette figure atteste le \emph{voisinage} du seuil, non une stabilité
d'ensemble, et la décroissance complète est mesurée au \S\ref{soc:sec:validation}.
\textit{Droite} : le nombre d'excès correspondant ($n = 91$) demeure au-dessus du
minimum de $30$ requis pour une estimation robuste.}
\label{soc:fig:diagnostic-gpd}
\end{figure}

\section{Hill contre maximum de vraisemblance : deux estimateurs confrontés}

% SOURCES-SCRIPTS: 47 63 67

L'estimateur de Hill \citep{Hill1975}, appliqué au même périmètre, ne corrobore pas l'indice de
queue par sa valeur : au seuil retenu ($k=86$) il vaut $1{,}32$, soit $120{,}8\,\%$ au-dessus du
maximum de vraisemblance. L'écart est pourtant exactement ce qu'un indice de $0{,}595$ produit en
échantillon fini : sur $2\,000$ échantillons simulés sous le modèle publié, les six valeurs
observées de Hill tombent dans l'intervalle simulé à $90\,\%$. Hill croît en outre avec $k$ quand
le balayage de seuil fait décroître $\hat\xi$, ce qui serait inversé si la queue valait $1{,}32$.
La calibration sur la PRC donne $\hat\xi = 1{,}03$, un indice qui ne porte pas sur le même
objet, la sévérité PRC étant une puissance du nombre d'enregistrements (\S\ref{sec:oprisk}). La
définition de l'estimateur, la table des simulations, le tracé et la validation croisée sont au
\S\ref{sec:ann-hill}.

\section{Validation et adéquation du socle}
\label{soc:sec:validation}

Avant d'agréger, on vérifie que les deux lois posées s'ajustent réellement aux données, faute
de quoi tout le reste serait bâti sur du sable (figure~\ref{soc:fig:validation}).

\textbf{Sévérité.} La GPD passe les diagnostics sans réserve. Les tracés quantile-quantile et
probabilité-probabilité s'alignent sur la diagonale ; la vie résiduelle moyenne est linéaire
au-delà du seuil, ce qui est la signature même d'une queue GPD et justifie le choix de $u$ ;
et $\hat\xi$ est stable dans le \emph{voisinage} du seuil retenu, à $0{,}5\,\%$ près jusqu'au
percentile $85$ des données courantes. Il ne l'est pas sur toute la plage :
en remontant le seuil du percentile $85$ au percentile $95$, l'indice décroît de $0{,}60$ à
$0{,}38$, six seuils comptant chacun au moins trente excès. Cette décroissance reste
\emph{entièrement} dans l'intervalle de confiance à $90\,\%$ déjà publié pour $\xi$,
$[0{,}30\,;0{,}83]$ : la sensibilité au seuil ne crée donc pas d'incertitude nouvelle, elle se
lit dans celle qui est déjà déclarée. En descendant au contraire sous le seuil retenu, $\hat\xi$
remonte jusqu'à $0{,}98$ et sort de cet intervalle par le haut, ce qui est le biais de seuil
attendu lorsqu'on ajuste une loi de queue à des pertes qui n'y appartiennent pas encore, et ce
qui justifie de ne pas descendre plus bas. Les tests formels
ne rejettent pas l'ajustement, avec une $p$-value obtenue par bootstrap paramétrique
les valeurs critiques tabulées ne valant pas lorsque les paramètres sont estimés :
Anderson-Darling $p=0{,}96$, Kolmogorov-Smirnov $p=0{,}89$. Ces deux tests portent sur la
\emph{famille} GPD, paramètres réajustés à chaque tirage. On teste séparément le couple
$(\hat\xi,\hat\sigma)$ \emph{effectivement publié}, paramètres imposés, seule façon d'attester
le chiffre plutôt que le choix de loi : Anderson-Darling $p=0{,}99$, Kolmogorov-Smirnov
$p=0{,}89$, sans rejet non plus. Le seuil retenu ne porte pas la conclusion : au percentile
$85$ des données plutôt qu'au seuil publié, $\hat\xi$ ne bouge que de $0{,}5\,\%$. Enfin,
la couverture réelle
de l'intervalle de confiance à $90\,\%$ de $\xi$, mesurée par simulation à $n$ fini, vaut
$86\,\%$. Le mot juste n'est pas \og{}proche du nominal\fg{} : un intervalle annoncé à
$90\,\%$ qui n'en couvre que $86$ est trop étroit, donc l'incertitude reportée sur
$\xi$ est sous-estimée et non l'inverse. Sous l'approximation normale, atteindre le nominal
demanderait d'élargir la demi-largeur d'environ $10\,\%$. L'écart est petit, mais son sens
compte, et il rejoint celui du taux de dépassement gelé au \S\ref{sec:pu-coherence} ainsi que
celui de la dérive d'échelle mesurée au \S\ref{soc:sec:backtest} : les trois seuls défauts
d'estimation du mémoire vont du même côté, celui de la sous-estimation, quand tous ses choix
\emph{posés} vont du côté prudent. C'est aussi ce qui justifie de rapporter les
intervalles bootstrap plutôt que les asymptotiques.

\textbf{Fréquence.} Sur les $14\,944$ cellules firme-année du panel financier, la binomiale
négative domine le Poisson sans ambiguïté : le test du rapport de vraisemblance donne
$p\approx 10^{-30}$. La surdispersion ($\Var/\E=1{,}19$) est donc un \emph{fait}, non un
artefact, et la NB épouse la queue des comptes là où le Poisson s'effondre.

\begin{figure}[htbp]\centering
\figover{J3_validation_adequation.png}
\caption{Validation du socle. (a-b) QQ-plot et PP-plot de la sévérité contre la GPD
ajustée. (c) vie résiduelle moyenne, linéaire au-delà de $u$. (d) stabilité
de $\hat\xi$ au seuil. (e) test d'Anderson-Darling, loi nulle par bootstrap
paramétrique et statistique observée ($p=0{,}96$). (f) adéquation de la fréquence :
la NB colle aux comptes firme-année, le Poisson rate la queue.}
% SOURCES-SCRIPTS: 47 93
\label{soc:fig:validation}
\end{figure}

\paragraph{Le seuil est une règle, pas un choix.} Une règle de stabilité écrite, appliquée à dix
percentiles candidats, retient $19{,}87$~M\euro{} contre $20{,}03$ publiés, avec un quantile
identique, et ce seuil est aussi le mieux ajusté du balayage. La règle la plus permissive
descendrait à $12{,}88$~M\euro{} et donnerait un quantile de sévérité supérieur de $16{,}0\,\%$ :
le seuil publié n'est pas celui qui maximise le chiffre du mémoire. Les deux règles et l'absence
de compromis entre biais et variance sur cette grandeur sont au \S\ref{sec:ann-seuil}.

\begin{cle}
\textbf{Ce que la validation établit, et ce qu'elle n'établit pas.} Les tests portent sur la
\emph{forme} des lois, pas sur le \emph{niveau} du capital : ils confirment que la GPD et la
NB sont les bons objets, mais laissent entière l'incertitude sur $\xi$ et donc sur la VaR (la
bande du chapitre~\ref{chap:resultats}). Adéquation n'est pas certitude : c'est précisément
pourquoi le SCR se lit en intervalle, jamais en point.
\end{cle}

%--------------------------------------------------------------

\section{Le socle hors échantillon : un backtest à origine glissante}
\label{soc:sec:backtest}

Les diagnostics qui précèdent partagent une limite, et elle est de principe plutôt que de
degré : ils interrogent l'ajustement sur les données mêmes qui l'ont produit. Un test
d'adéquation demande si une loi décrit son échantillon, jamais si elle prédit une période
qu'elle n'a pas vue. C'est pourtant la seconde question qu'un lecteur adresse à un modèle de
capital, et cette section y répond.

Le protocole est celui de l'origine glissante. Pour chaque année d'origine, les deux lois sont
ajustées sur les seules observations antérieures, la seule année suivante est notée, puis
l'origine avance d'un an. La fenêtre retenue va de $2004$ à $2025$, soit $539$ incidents, ce qui
donne douze années notées hors échantillon. Deux précautions en conditionnent la lecture. La
fenêtre est choisie sur la donnée et non par commodité : avant $2004$ la base porte de un à neuf
incidents par an contre treize à trente-neuf ensuite, et l'année $2026$ n'en porte que cinq, si
bien qu'une fenêtre plus large fabriquerait un faux régime calme au début et un faux
effondrement à la fin. Et le seuil est ré-estimé sur chaque échantillon d'apprentissage
par la règle du percentile $85$, plutôt que lu dans la calibration figée : un seuil calculé sur
toute la période ferait passer l'information de test dans l'apprentissage, et le test cesserait
de mesurer quoi que ce soit.

\paragraph{Les résultats.} La loi de fréquence est validée : un intervalle de prévision annoncé à
$90\,\%$ en couvre $91{,}7\,\%$ sous la binomiale négative contre $75{,}0\,\%$ sous le Poisson,
et le log-score, qui pénalise la largeur, désigne le même gagnant. La forme de la queue survit au
test de probabilité intégrale ($p=0{,}4087$). Les quantiles de sévérité sont en revanche dépassés
trois fois trop souvent, parce que l'échelle dérive : $4{,}00\,\%$ par an dans la queue, dont la
moitié d'inflation monétaire, les montants de la base étant nominaux, soit au moins
$+24{,}7\,\%$ sur le quantile de sévérité. Cette dérive ne gonfle pas la thèse : le facteur entre
états passe de $3{,}344$ à $3{,}124$ sous la sévérité dérivée, et l'écart en euros ne bouge que
de $-1{,}9\,\%$, par compensation entre forme et échelle. Sur la charge annuelle agrégée le
verdict est négatif, pour la même cause, le désalignement n'apparaissant que sur les dernières
années notées. L'indépendance entre fréquence et sévérité tient sur les excès. Le détail de ces
tests est au \S\ref{sec:ann-backtest}.
% SOURCES-SCRIPTS: 89 90 91 102

\begin{attention}
\textbf{Ce que ce dispositif ne peut pas tester, dit en années plutôt qu'en précaution de style.}
Détecter un taux de dépassement \emph{double} du nominal avec une puissance de $80\,\%$
demanderait $1\,811$ années au niveau de $99{,}5\,\%$, $905$ à $99\,\%$, et $78$ même à
$90\,\%$. Le quantile qui porte le capital n'est donc backtestable sur aucun historique
de risque opérationnel existant, et sur douze années le test ne renseigne que le centre et le
corps de la loi de l'agrégat. Le niveau de capital reste adossé à l'ajustement paramétrique,
dont la queue est corroborée par trois chemins indépendants et bornée par l'intervalle de
$\xi$ ; aucun de ces chemins n'est un backtest du quantile, et le présenter comme tel serait
faux. Un backtest qui ne déclare pas sa puissance laisse croire qu'une absence de rejet vaut
validation.
\end{attention}

\begin{cle}
\textbf{Ce que le backtest établit, et il y a un seul défaut derrière trois symptômes.} La loi de
fréquence est validée hors échantillon, sur deux critères dont un que la largeur n'achète pas. La
forme de la queue survit au test de probabilité intégrale. L'indépendance entre fréquence et
sévérité tient sur la population que le modèle tire. Un seul mécanisme est en défaut,
l'échelle de la sévérité, qui dérive de $4{,}00\,\%$ par an dans la queue, et les trois symptômes
qui restent en découlent : les quantiles de sévérité dépassés trois fois trop souvent, la charge
agrégée rejetée, et un quantile publié anti-conservateur d'au moins $24{,}7\,\%$. La valeur reste
gelée et l'écart est inscrit au registre des limites (\S\ref{sec:pu-coherence}), du même
traitement que le taux de dépassement. Sur la thèse, la formule à retenir est
l'invariance d'échelle et non l'invariance à la dérive : le facteur entre états ne bouge pas
d'une variation d'échelle, mais il baisse de $6{,}6\,\%$ sous la sévérité dérivée, dont la forme
est plus légère.
% SOURCES-SCRIPTS: 89 90 91
\end{cle}

%--------------------------------------------------------------

% =====================================================================
